\documentclass[11pt]{article}
\usepackage[margin=1in]{geometry}
\usepackage{amsmath,amssymb}
\usepackage{graphicx}
\usepackage{subcaption}
\usepackage{titlesec}
\usepackage{enumitem}
\usepackage{comment}
\usepackage[
  backend=biber,
  style=numeric,
  sorting=none
]{biblatex}
\addbibresource{main_fixed.bib}
\usepackage{wrapfig}
\usepackage{float}
\usepackage{overpic}
\usepackage{tabularx}
\usepackage{array}
\usepackage{xcolor}

\titleformat{\section}{\large\bfseries}{\thesection}{0.5em}{}
\titleformat{\subsection}{\normalsize\bfseries}{\thesubsection}{0.5em}{}

\begin{document}
\begin{center}
  {\Large \textbf{Active Anchoring at Solid Walls in Simulated Active Nematics}}\\[2ex]
  Michael Fang \\
  \today
\end{center}

\noindent
{\small
\begin{minipage}{\textwidth}
\centering

\begin{minipage}[t]{0.44\textwidth}\vspace{0pt}
Primary advisor: Professor William Bialek\\
Second reader: Professor Andrej Košmrlj\\
\begin{tabular}[t]{@{}l@{\ }l@{}}
External advisors: & Professor Julia Yeomans\\
                  & Dr.\ Jan Rozman
\end{tabular}
\end{minipage}
\hspace{0.04\textwidth}%
\begin{minipage}[t]{0.40\textwidth}\vspace{0pt}
\raggedleft
This paper represents my own work in accordance with University regulations.\\
/s/\ Michael Fang
\end{minipage}

\end{minipage}
}
\vspace{-0.2cm}

\begin{abstract}
Active nematics exhibit distinct orientation patterns near boundaries, where activity generates shear flows that select preferred director alignments. We study this active anchoring phenomenon using continuum theory and hybrid lattice-Boltzmann simulations implemented in Ludwig. Starting from the Stokes and Leslie--Ericksen framework for nemato-hydrodynamics, we first review how active stresses, flow-alignment, and elasticity affect stable orientations at active-passive boundaries. We then analytically determine how active anchoring occurs at solid walls and test these predictions by simulating two-dimensional active nematic systems. Our simulated results reproduce the theoretical selection of planar anchoring in extensile systems and homeotropic anchoring in contractile systems, and show how increasing the flow-alignment parameter shifts the stable director angle at solid boundaries from active anchoring towards a well-defined Leslie angle relative to the flow. By measuring the depth of active anchoring away from the boundary walls, we also confirm that the predicted penetration depth of active anchoring is proportional to the characteristic active length scale $l_\text{act} \sim \sqrt{\kappa / |\zeta |}$. These results provide a direct numerical test of solid-wall active anchoring in confined active nematics.

\end{abstract}

\tableofcontents
\clearpage

\section{Introduction}

All living systems stay alive by constantly drawing energy from their environment and using it to do work. These systems can be classified as active matter, which describes non-equilibrium systems that draw energy from their surroundings to move, push on their surroundings, and reorganize themselves, ranging from biological cells and microswimmers to flocks of fish or birds~\cite{ramaswamy2010mechanics}. At small scales, these processes can be observed in many ways, from bacteria swimming by rotating their flagella to cells crawling by extending their actin networks~\cite{yeomans2025active,saw2017topological}. Because each unit continuously draws energy from its surroundings, these systems are able to sustain flows and structures that would otherwise die out~\cite{doostmohammadi2018active}. Understanding active matter therefore helps reveal how living systems generate organization, motion, and behavior out of equilibrium. 

Within the broad family of active matter, it is useful to think about the shape and symmetry of the constituent particles, as they are often non-spherical. In fact, many systems are made up of elongated rod-like particles with head-tail symmetry. These features identify the particles as resembling a state of matter known as nematic liquid crystals, which at high concentrations exhibit long-range orientational order, but flow under application of shear stress~\cite{degennes1993physics}. This means that the rods are able to move around relative to each other, but unlike traditional liquids, nematics tend to align with their neighbors along their axes of elongation. Because the rods are head-tail symmetric, flipping every rod end-for-end has no effect on the overall state. The nature of nematic particles to align orientationally is reminiscent of crystal structures, hence their label as liquid crystals. 

\begin{figure}[H]
  \centering
  \includegraphics[width=0.58\textwidth]{gemini_nematics2.png}
  \caption{Schematic representation of ordering in (a) crystalline solid, (b) isotropic liquid, and (c) nematic liquid crystal phases.}
  \label{fig:gemini_nematics}
\end{figure}

The intersection between active matter and nematic liquid crystals is the field of active nematics, where elongated self-driven units resembling nematic particles continuously generate internal stresses and spontaneous flows. At large enough density, the rod-shaped particles locally share a common axis, but because each unit is powered, the state of the system is constantly changing. Some examples of active nematic systems include layers of bacteria at high concentration and crowded sheets of elongated cells, like epithelial monolayers and tissues~\cite{doostmohammadi2018active, nejad2024stress}. Thus, constructing a framework to model and understand active nematics allows us to not only identify why certain behaviors come about (from activity vs shape), but to also understand how systems like tissues healing and microswimmers steering and sorting work. 

A key characteristic of active nematics is their activity, which is represented by an activity parameter $\zeta$. The magnitude of $\zeta$ determines the strength of the forces and shear exerted by the active nematics on their surroundings, while the sign distinguishes extensile systems ($\zeta>0$, which push fluid outward along their axis and draw it in from the sides) from contractile systems ($\zeta<0$, which do the opposite). The forces and flows associated with extensile and contractile particles are shown in Fig.~\ref{fig:gemini_nematic_forces} below. 

\begin{figure}[H]
  \centering
  \includegraphics[width=0.65\textwidth]{gemini_nematic_forces2.png}
  \caption{Active forces (in black) and resulting flows (in blue) for (a) extensile ($\zeta>0$) and (b) contractile ($\zeta<0$) active nematic particles.}
  \label{fig:gemini_nematic_forces}
\end{figure}

One particular aspect of active nematics worth studying is interactions at walls and interfaces. In real life, active materials like bacterial films and cell layers almost always live in channels, droplets, or on surfaces. Near boundaries, active nematic units often prefer certain orientations, with their continual pushing or pulling driving flows along surfaces~\cite{blow2014biphasic,bhattacharyya2023phase}. The organization of active nematics at active-passive interfaces exhibits ``active anchoring," which motivates the boundary-focused analysis that will follow in this paper.

\begin{figure}[H]
  \centering
  \includegraphics[width=0.45\textwidth]{active_anchoring_interface2.png}
  \caption{Extensile ($\zeta>0$) active nematics at an interface generate a tangential force that aligns them parallel to the boundary~\cite{bhattacharyya2023phase,Bhattacharyya2024ActiveNematicsMechanobiology}.}
  \label{fig:active_anchoring_interface}
\end{figure}

While the mechanism of active anchoring is well-understood at active-passive interfaces, many experimental and biological realizations are confined by solid boundaries rather than passive isotropic phases. The main mechanism by which typical active anchoring comes about, which is spatial gradients in nematic order near active-passive interfaces, does not apply in the same way in solid-wall systems as the nematic order does not necessarily vanish at the boundary. Thus, it is unclear whether the same active anchoring selection found at active-passive interfaces should also emerge at solid walls, or how it competes with flow-alignment effects in confined geometries. 

In this work, we address this question by studying a confined active nematic in a channel with solid, no-slip walls and no imposed surface anchoring. We analyze how activity and flow-alignment compete to select wall orientations using a continuum description based on Stokes flow and Leslie--Ericksen director dynamics, and we test these predictions with lattice-Boltzmann simulations.

\section{Theory of Active Nematics}

In this section, we cover the continuum framework needed to analyze active anchoring at active-passive interfaces, and to interpret the solid-wall results that follow. We first introduce the hydrodynamic and orientational fields used to describe active nematics and set up the active stress that couples orientation gradients to flow. The key mechanism by which we model active nematics is through the coupled progression of a flow field $\mathbf{u}$ and an orientation field $\mathbf{Q}$. We then analyze the near-interfacial force balance in the Stokes regime to obtain the active shear generated at interfaces, and combine this with the Leslie--Ericksen equation to determine how the director responds to shear and selects a stable anchoring angle. We review in-depth active anchoring at active-passive interfaces because the same stress-gradient mechanism will reappear for solid walls in a confined channel. Throughout this report, we will work in two dimensions and in the low-Reynolds-number (Stokes) limit, which is typical for dense active nematics~\cite{yeomans2025active,doostmohammadi2018active}. 

\subsection{Fields and order parameters}
\label{sec:fields_and_order_parameters}

We will operate in the continuum description of active nematics, which builds on liquid crystal hydrodynamics, extended to include active stress in its modified stress tensor. To start, we replace the many individual swimmers in a system with a few smooth fields obtained by coarse-graining over small neighborhoods~\cite{doostmohammadi2018active, simha2002hydrodynamic}. Coarse-graining simply means we are averaging over enough particles that rapid microscopic fluctuations cancel, while slow spatial variations remain. This is useful in confined geometries because walls influence the system over distances large compared with a single bacterium, so predictions about alignment and flow should depend on local averages rather than on the exact positions and orientations of every particle. 

The first field we use is the velocity field $\mathbf{u}(\mathbf{x},t)$, which describes the collective motion of the active nematic fluid. We assume the active system is incompressible, with constant density. By conservation of mass, the velocity field must be divergence-free, so 
\begin{equation}
    \nabla\!\cdot\!\mathbf{u}=0.
\end{equation}
By force balance, we also have that
\begin{equation}
\rho\big(\partial_t+\mathbf{u}\!\cdot\!\nabla\big)\mathbf{u}=\nabla\!\cdot\!\boldsymbol{\Pi},
\end{equation}
where the left-hand side is the convected derivative representing the total rate of change of momentum of a fluid element in flow, and the right side is the derivative of the stress tensor, which contains the forces that conserve the total momentum of the fluid. In the low-Re regime relevant here, we can drop the inertial (left-hand) side and work with $\nabla\cdot\boldsymbol{\Pi}=0$. The components of the stress tensor are elaborated on in the following section. 

The second field we use to describe an active nematic system is the orientation field  (or director field) $\mathbf{n}(\mathbf{x},t)$, which describes the average axis of elongation in position and time. Because each rod-shaped particle is head-tail symmetric, the only physically meaningful object is its axis, motivating the use of a headless director $\mathbf{n}(\mathbf{x},t)$ with $\mathbf{n} = -\mathbf{n}$ and $|\mathbf{n}|=1$~\cite{degennes1993physics}. Thus, while a normal vector would be invariant under $\theta \rightarrow \theta + 2\pi$, our director field vector $\textbf{n}$ is invariant under $\theta \rightarrow \theta + \pi$.

For our coarse-grained model, simply knowing the average axis of orientation is not sufficient. We can see this by imagining two systems with the same average orientation direction, but one with strong local alignment and the other barely ordered, resulting in very different stresses and behavior. This motivates the introduction of a scalar measure of how sharply local orientations cluster about the average orientation direction. 

In two dimensions, letting $\theta=0$ represent an orientation aligned along the x-axis, we can denote $f(\theta)$ to be the local probability density of in–plane orientation angles $\theta$ in a small box around $\mathbf{x}$. Then, $f(\theta)=f(\theta+\pi)$ captures head-tail symmetry, and the director angle is the location where $f(\theta)$ is most concentrated. We can then define the nematic order parameter $S(\mathbf{x},t)\in[0,1]$ in two dimensions as 
\begin{equation}
S = \int_{0}^{\pi} f(\theta)\big(2\cos^{2}\theta-1\big)d\theta
=\int_{0}^{\pi} f(\theta)\cos 2\theta d\theta,
\end{equation}
which yields $S=1$ for perfect alignment and $S=0$ for uniformly random orientations~\cite{degennes1993physics}. This value respects head-tail symmetry (with period $\pi$ in angle), increases when $f(\theta)$ becomes more peaked, and disappears when orientations are uniformly random. 

We now have a director field $\textbf{n}(\textbf{x},t)$ that describes average orientation and a scalar field $S(\textbf{x},t)$ that describes nematic order. While these fields are intuitive, it is inefficient to write hydrodynamics purely in terms of a headless axis and a scalar. To fix this, we can define a rotationally covariant object that packs information on both the orientation direction and order into a single field. This finally then motivates the use of a symmetric, traceless nematic order tensor 
\begin{equation}
\mathbf{Q}(\mathbf{x},t)=\frac{d}{d-1}S(\mathbf{x},t)\!\left(\mathbf{n}\mathbf{n}-\frac{\mathbf{I}}{d}\right),
\end{equation}
in $d$ spatial dimensions~\cite{degennes1993physics, beris1994thermodynamics, Bhattacharyya2024ActiveNematicsMechanobiology}. $\mathbf{Q}$ is unchanged by $\mathbf{n}\to-\mathbf{n}$ and has principal eigenvector $\mathbf{n}$ with eigenvalues proportional to $S$ capturing the nematic order, while the final term proportional to $\mathbf{I}$ is subtracted to keep $\mathbf{Q}$ traceless, normalizing $\mathbf{Q}=0$ to correspond to no nematic order. In two dimensions, we can write $\mathbf{n}=(\cos\theta,\sin\theta)$, which yields $$Q_{xx}=-Q_{yy}=S\cos 2\theta,\qquad Q_{xy}=Q_{yx}=S\sin 2\theta,$$ and therefore
\begin{equation}
    \mathbf{Q} = S\begin{pmatrix}
    \cos 2\theta & \sin 2\theta\\
    \sin 2\theta & -\cos 2\theta
    \end{pmatrix}
\end{equation}
which is indeed symmetric and traceless. Now that we have defined the orientation field $\mathbf{n},$ the velocity field $\mathbf{u},$ the order parameter S, and the nematic order tensor $\mathbf{Q}$, let us move on to the hydrodynamic equations that couple $\mathbf{u}$ and $\mathbf{Q}$. 

\subsection{Nemato-hydrodynamics}
\label{sec:nemato-hydrodynamics}

The orientational field and the fluid flow must evolve together because each one changes the other. The movement of the surrounding fluid rotates and stretches the local alignment of rods, while spatial variations of alignment, as well as activity, in turn push on the liquid. Even though forces balance, these interactions generate dynamics because $\mathbf{u}$ responds instantaneously to $\mathbf{Q}$, but $\mathbf{Q}$ can take time to respond to $\mathbf{u}$. The goal of this section is to establish a system of equations that capture this two-way coupling.

In dense active nematic systems, typical speeds $v$ and lengths $L$ give Reynolds number $\mathrm{Re}\sim \rho vL/\eta\ll 1$ so inertial effects are negligible compared with viscous forces~\cite{simha2002hydrodynamic}. Thus, at each instant, the fluid is effectively incompressible and the flow adjusts to balance internal stresses to give:
\begin{equation}
\nabla\!\cdot\!\mathbf{u}=0, \qquad \nabla\!\cdot\!\boldsymbol{\Pi} = \mathbf{0},
\label{eq:navier-stokes}
\end{equation}

To describe how a local axis responds to deformation, we split the velocity gradient into its symmetric and antisymmetric parts~\cite{doostmohammadi2018active, beris1994thermodynamics, Bhattacharyya2024ActiveNematicsMechanobiology},
\begin{equation}
    \mathbf{E}\equiv\tfrac12\!\big(\nabla\mathbf{u}+(\nabla\mathbf{u})^{T}\big),\qquad \boldsymbol{\Omega}\equiv\tfrac12\!\big(\nabla\mathbf{u}-(\nabla\mathbf{u})^{T}\big).
\end{equation}

The strain rate tensor $\mathbf{E}$ measures the local rate at which nearby fluid elements stretch and shear, while the vorticity tensor $\boldsymbol{\Omega}$ measures the local angular velocity for rigid rotation. 

The first contribution to the total stress tensor comes from viscous dissipation. When neighboring layers of fluid move relative to one another, they experience drag forces that oppose this motion and transfer momentum from faster to slower regions. These internal frictional forces can be represented by a viscous stress 
\begin{equation}
    \boldsymbol{\Pi}_{\mathrm{visc}} = 2\eta\mathbf{E},
\end{equation}
where $\mathbf{E}$ is the strain rate tensor defined above and $\eta$ is the shear viscosity of the fluid, which represents the relative significance between this deformation rate and the resulting viscous stress~\cite{degennes1993physics}. 

The next contribution to the total stress tensor is the active stress, which is obtained through the coarse-grained effect of many microscopic force dipoles~\cite{simha2002hydrodynamic}. The explicit derivation can be found in Appendix A, but the final coarse-grained result gives our second stress term 
\begin{equation}
    \boldsymbol{\Pi}_{\mathrm{act}} = -\zeta \textbf{Q},
\end{equation}
where $\zeta$ is the activity coefficient that determines the strength of the activity in a given region. In the scope of this paper, we take $\zeta$ to be constant with approximately uniform particle strength and density.

A further contribution to the total stress tensor is the elastic, or ``backflow" term. Spatial variations of the orientation field $\mathbf{Q}$ cost free energy, as neighboring rods prefer to align with each other instead of twisting or splaying~\cite{degennes1993physics}. As $\mathbf{Q}$ changes in space and time, these variations then generate restoring torques and exert stresses on the surrounding fluid through backflow. This effect is represented by an elastic stress $\boldsymbol{\Pi}_{\mathrm{el}}$, which can be derived from a standard Landau--de Gennes free energy functional for $\mathbf{Q}$~\cite{degennes1993physics, Bhattacharyya2024ActiveNematicsMechanobiology, beris1994thermodynamics}:
\begin{equation}
\mathcal{F}[\mathbf{Q}] \;=\; \int d^{2}r \left[
C\left(S_{\mathrm{nem}}^{2}-\frac{1}{2}\,\mathrm{Tr}\!\left(\mathbf{Q}^{2}\right)\right)^{2}
+\frac{\kappa}{2}\,(\partial_\gamma Q_{\alpha\beta})(\partial_\gamma Q_{\alpha\beta})
\right].
\label{eq:free_energy}
\end{equation}

The bulk term is a Landau--de Gennes potential minimized at $\mathrm{Tr}(\mathbf{Q}^2)=2S_{\mathrm{nem}}^2$, which fixes the preferred magnitude of nematic order, while the gradient term penalizes spatial distortions of $Q$ and introduces the elastic constant $\kappa$ that controls how quickly wall--induced variations relax into the bulk. This same functional defines the molecular field $\mathbf{H}$ that appears in the Beris--Edwards dynamics through a functional derivative (as written in Eq.~\eqref{eq:Beris-Edwards} below).

The last contribution is the isotropic pressure term $-p \mathbf{I}$, which enforces incompressibility, \mbox{($\nabla\cdot\mathbf{u}=0$)}. While pressure contributes to normal force balance, it does not enter the tangential momentum equation in the 1D shear geometry $u=(u_x(y),0)$ with $\partial_x(\cdot)=0$, as $(\nabla\cdot(-p\mathbf{I}))_x = -\partial_x p = 0$.
Finally, we can fully construct the total stress tensor for active nematic hydrodynamics:
$$\boldsymbol{\Pi} = \boldsymbol{\Pi}_{\text{visc}} + \boldsymbol{\Pi}_{\text{act}}+\boldsymbol{\Pi}_{\text{el}}-p\mathbf{I}$$
\begin{equation}
    = 2\eta\mathbf{E} - \zeta\mathbf{Q} + \boldsymbol{\Pi}_{\mathrm{el}} -p\mathbf{I}
\end{equation}
In the boundary problems considered in this paper, we are concerned with near-interfacial/near-wall shear and alignment, so in the tangential momentum balance we can drop the pressure term, as it does not affect shear. We also drop the elastic stress term in the tangential direction because it is negligible relative to the viscous and active terms. Eq. \eqref{eq:navier-stokes} then gives us: 
\begin{equation}
\nabla\!\cdot\!\big(2\eta\mathbf{E}-\zeta\mathbf{Q}\big)=0,
\qquad
\nabla\!\cdot\!\mathbf{u}=0.
\end{equation}

To describe the evolution of the orientational field itself, we note that $\mathbf{Q}$ is not fixed but evolves under both elastic relaxation and flow-induced reorientation. In general, the dynamics of a nematic field are governed by the Beris–Edwards equation~\cite{beris1994thermodynamics}, which couples the time evolution of $\mathbf{Q}$ to the local velocity gradients:
\begin{equation}
\partial_t \mathbf{Q} + \mathbf{u}\!\cdot\!\nabla \mathbf{Q} - \mathcal{W}(\mathbf{Q},\nabla\mathbf{u})
= \Gamma\mathbf{H}.
\label{eq:Beris-Edwards}
\end{equation}
The term 
\begin{equation}
    \mathcal{W}(\mathbf{Q},\nabla\mathbf{u}) = (\lambda \mathbf{E} + \boldsymbol{\Omega}) \cdot \left(\mathbf{Q} + \frac{\mathbf{I}}{d}\right) + \left(\mathbf{Q} + \frac{\mathbf{I}}{d}\right) \cdot (\lambda \mathbf{E} - \boldsymbol{\Omega}) - 2\lambda \left(\mathbf{Q} + \frac{\mathbf{I}}{d}\right)\,(\mathbf{Q}:\mathbf{E})
    \label{eq:co--rotation}
\end{equation}
is a co--rotational term that describes how $\mathbf{Q}$ is advected, rotated, and stretched in the flow.

$\Gamma$ is a collective rotational diffusivity controlling the relative importance of the free energy in the dynamics of Q (effectively controlling how quickly the order field relaxes), and $\mathbf{H}$ is the molecular field derived from the variation of the Landau--de Gennes free energy $\mathcal{F}$ that drives $\mathbf{Q}$ toward equilibrium. 
\begin{equation}
\mathbf{H}
= -\left(\frac{\delta \mathcal{F}}{\delta \mathbf{Q}}
- \frac{\mathbf{I}}{d}\,\mathrm{Tr}\!\left(\frac{\delta \mathcal{F}}{\delta \mathbf{Q}}\right)\right).
\end{equation}

In the scope of this paper, we are focusing on steady near-boundary configurations where $\mathbf{n}$ can be approximated as locally uniform. Thus, we can get by without explicitly writing out or solving the Beris-Edwards equation. When we later solve for active anchoring at solid walls in the results section, we will use the Leslie--Ericksen equation for director dynamics, which is obtained as a limit of Beris--Edwards for approximately constant $S$ and uniaxial $\mathbf{Q}$~\cite{degennes1993physics, leslie1968some}. In this limit, we use a local torque balance neglecting the molecular-field/elastic term that would require solving the full $\mathbf{Q}$ PDE. These elastic effects then enter separately when we discuss the penetration depth of the wall--induced distortion in Sec.~\ref{sec:penetration_depth}.

\subsection{Active anchoring at active-passive interfaces}
\label{sec:active_anchoring_interfaces}

In passive liquid crystals, boundaries influence orientation through surface interactions, as chemical treatments or roughness favor particular directions of alignment at an interface or wall. For passive nematics this is known as anchoring. In active nematics, because activity generates stresses wherever the orientational field varies, any interface where the order parameter changes rapidly can experience active forces that, by themselves, tend to align the material in preferred directions~\cite{blow2014biphasic,Bhattacharyya2024ActiveNematicsMechanobiology}. This boundary-driven selection of orientation caused by active stresses is called active anchoring. 

In this section, we briefly re-derive the existing result for active anchoring at a nematic--isotropic interface. We start with the active contribution to the stress tensor, $$\boldsymbol{\Pi}^{\mathrm{act}} = -\zeta\mathbf{Q},$$
where $\zeta$ is the activity coefficient and $\mathbf{Q}$ is the nematic order tensor. Because the momentum equation in the Stokes regime involves the divergence of $\boldsymbol{\Pi}$, the corresponding active force density (force per unit volume) is
\begin{equation}
\mathbf{f}^{\mathrm{act}} = \nabla\!\cdot\!\boldsymbol{\Pi}^{\mathrm{act}} = -\zeta\nabla\!\cdot\!\mathbf{Q}.
\end{equation}

In the bulk, where $\mathbf{Q}$ is approximately uniform, these forces are small, but near boundaries or interfaces where $\mathbf{Q}$ can vary sharply over short distances, the active forces can be large and organized. At an interface between an active nematic and an isotropic phase, let $\mathbf{m}$ be the unit normal to the interface, pointing from nematic into isotropic, and let $\mathbf{l}$ be a unit tangent lying in the interface plane. At any point along the interface, let $\mathbf{n}$ be the orientation direction of the nematic, which points at an angle $\theta$ with respect to $\mathbf{l}$.\footnote{Credit to Saraswat Bhattacharyya~\cite{Bhattacharyya2024ActiveNematicsMechanobiology} for the setup in this derivation.} A diagram of the nematics at an active-passive interface is shown in Fig.~\ref{fig:gemini_nematic_boundary} below.

\begin{figure}[H]
  \centering
  \includegraphics[width=0.38\textwidth]{gemini_nematic_boundary8.png}
  \caption{Diagram for active nematic at a boundary.}
  \label{fig:gemini_nematic_boundary}
\end{figure}

We assume that near the interface, the order parameter $S$ varies only in the normal direction, while the director $\mathbf{n}$ is approximately uniform. From here, we can solve for the active forces at the boundary. We begin by decomposing the director into components along $\mathbf{m}$ and $\mathbf{l}$, 
\begin{equation}
    \mathbf{n} = (\mathbf{n}\!\cdot\!\mathbf{m})\mathbf{m} + (\mathbf{n}\!\cdot\!\mathbf{l})\mathbf{l}.
\end{equation}

The angle $\theta$ between $\mathbf{n}$ and the interface follows $\cos\theta = \mathbf{n}\!\cdot\!\mathbf{l}$ and $\sin\theta = \mathbf{n}\!\cdot\!\mathbf{m}$. Let us define $n_x$ and $n_y$ to be $n_x = \cos\theta$ and $n_y = \sin\theta$, so that we have $\mathbf{n} = (\cos\theta,\sin\theta) = (n_x,n_y)$, with $n_x^2+n_y^2=1$. In two dimensions, we can then write the active stress as
$$\boldsymbol{\Pi}_{\mathrm{act}} = -\zeta\mathbf{Q}
= -\zeta\Big[2S\big({\mathbf n}{\mathbf n} - \tfrac{\mathbf I}{2}\big)\Big], \quad{\mathbf n}=(n_x,n_y),\quad n_x^2+n_y^2=1.$$

Then, the active force is
\begin{equation}
    \mathbf f_{\mathrm{act}} = \nabla\!\cdot\boldsymbol{\Pi}_{\mathrm{act}} = -\zeta\nabla\!\cdot\!\Big(2S({\mathbf n}{\mathbf n}-\tfrac{\mathbf I}{2})\Big) = -2\zeta \nabla\!\cdot \Big( S \Big({\mathbf n}{\mathbf n}-\tfrac{\mathbf I}{2}\Big)\Big)
\end{equation}

Using $\nabla\!\cdot(2S\mathbf A)=2(\nabla S)\!\cdot\!\mathbf A + 2S(\nabla\!\cdot\mathbf A)$ and $\nabla\!\cdot\mathbf A=\mathbf 0$ for constant $\mathbf n$, we get
\begin{equation}
    \mathbf f_{\mathrm{act}} = -2\zeta\Big[(\nabla S)\!\cdot\!\big({\mathbf n}{\mathbf n}-\tfrac{1}{2}\mathbf I\big)\Big]
\end{equation}

Because we assume S only varies in the wall-normal direction $y$, we have that
\begin{equation}
    \nabla S = \begin{pmatrix} 0\\[2pt] S_y \end{pmatrix}, \qquad {\mathbf n}{\mathbf n}-\tfrac{1}{2}\mathbf I = \begin{pmatrix} n_x^2-\tfrac{1}{2} & n_x n_y\\[2pt] n_x n_y & n_y^2-\tfrac{1}{2} \end{pmatrix}
\end{equation}

Plugging back into our force equation, we then have
\begin{align}
\mathbf f_{\mathrm{act}}
&= -2\zeta\left[\begin{pmatrix} 0\\[2pt] S_y \end{pmatrix}
\begin{pmatrix} n_x^2-\tfrac{1}{2} & n_x n_y\\[2pt] n_x n_y & n_y^2-\tfrac{1}{2} \end{pmatrix}\right] \notag \\
&= -\zeta |S_y| \begin{pmatrix} 2n_x n_y\\[2pt] 2n_y^2 - 1 \end{pmatrix} \notag \\
&= -\zeta |\nabla S|\left[ 2n_y \begin{pmatrix} n_x\\[2pt] n_y \end{pmatrix}
- \begin{pmatrix} 0\\[2pt] 1 \end{pmatrix}\right] \notag \\
&= -\zeta |\nabla S| \left[ 2 (\mathbf{m}\cdot\mathbf{n})\,\mathbf{n} - \mathbf{m}\right].
\label{eq:fact_final}
\end{align}

Splitting this into its tangential and normal components, we have:
\begin{equation}
    \mathbf{f}_{\mathrm{act}}\!\cdot\!\mathbf{l} = -\,\zeta\,|\nabla S|\,\Big[\,2(\mathbf{m}\!\cdot\!\mathbf{n})(\mathbf{l}\!\cdot\!\mathbf{n})\Big] = -\,\zeta\,|\nabla S|\,\Big[\,2\sin \theta \cos \theta\Big] = -\,\zeta\,|\nabla S|\,\sin 2\theta, \end{equation}
\begin{equation}
    \mathbf{f}^{\mathrm{act}}\!\cdot\!\mathbf{m} = -\,\zeta\,|\nabla S|\,\Big[\,2(\mathbf{m}\!\cdot\!\mathbf{n})^{2}-1\Big] = -\,\zeta\,|\nabla S|\,\big(2\sin^{2}\theta-1\big) = -\,\zeta\,|\nabla S|\,\cos 2\theta.
\end{equation}

The tangential component $-\zeta|\nabla S|\sin 2\theta\,\mathbf{l}$ is the one that drives a near-wall shear flow that rotates the director towards its anchored configuration. The zeroes of $\sin 2\theta$ are at $\theta=0^\circ$ (director parallel to the interface, called ``planar" anchoring) and $\theta=90^\circ$ (director perpendicular to the interface, called “homeotropic” anchoring). The sign of $\zeta$ determines which of these is stable: for extensile activity $(\zeta>0)$ the induced shear rotates $\mathbf{n}$ towards $\theta=0^\circ$ (planar anchoring), while for contractile activity $(\zeta<0)$ the directors are rotated towards $\theta=90^\circ$ (homeotropic anchoring)~\cite{doostmohammadi2018active}. In either case, small deviations from the preferred orientation reverse the sign of the tangential force and the restoring shear, returning the director to the aligned state. 

Finally, it is important to note that this ``active anchoring" can occur even in the absence of any explicit surface energy favoring a particular orientation. While in passive nematics, anchoring is generally determined by minimizing a surface free energy term~\cite{degennes1993physics}, active anchoring arises from active stresses that generate forces wherever the nematic order varies spatially, following $\mathbf{f}^{\mathrm{act}}=\nabla\cdot\Pi^{\mathrm{act}}=-\zeta\,\nabla\cdot\mathbf{Q}$. In confined geometries, these effects are often both present and can compete to determine the preferred near-wall orientation of active nematics at interfaces and boundaries.

While in a solid-wall box or channel there is no active-passive interface where $S$ drops to zero, analogous wall-normal gradients in $\mathbf{Q}$ can arise in a near-wall boundary layer through variations in $S$, the director angle, or both. Then, because the active force density is $f^{\rm act}=\nabla\cdot\Pi^{\rm act}=-\zeta\,\nabla\cdot Q$, these wall-normal gradients generate the same shear-driven active anchoring mechanism as at active-passive interfaces. We will derive and test this result in later sections. 

\section{Simulation Methods}

\subsection{Hybrid lattice--Boltzmann framework}
\label{sec:LB-framework}

We simulate active nematic systems numerically by simultaneously solving the flow field $\mathbf{u}$ and orientation field $\mathbf{Q}$ using a hybrid lattice-Boltzmann method~\cite{marenduzzo2007steady}. This method combines a modified lattice-Boltzmann solver for the fluid velocity field with a finite difference solver for the nematic order tensor $\mathbf{Q}$ on the same Cartesian grid~\cite{Krueger2016LatticeBoltzmann}. All simulations are performed using Ludwig, an open--source LB code designed for simulating complex fluid systems that includes built--in modules for Beris–Edwards nematodynamics and active stresses.

For the flow field, we solve the Navier-Stokes and continuity equations using the lattice-Boltzmann algorithm with a BGK collision operator and additional source terms representing the active and elastic stresses~\cite{Krueger2016LatticeBoltzmann}. The simulated system is isothermal instead of exactly incompressible, but density fluctuations are rapidly damped by sound waves that propagate two orders of magnitude faster than any orientational or hydrodynamic timescale~\cite{KusumaatmajaYeomans2010LBWettingDrop}. As a result, the system is effectively incompressible with an error $\sim 0.01\%$. Finally, while inertial terms are present in the lattice-Boltzmann formulation, all simulations are run with Reynolds number Re $\approx 0.003$, so viscous stresses dominate and the dynamics remain in the low-Re regime relevant for active nematics. 

Simultaneously with the fluid solver, the nematic order tensor $\mathbf{Q}$ evolves according to the Beris-Edwards equation (Eq. \eqref{eq:Beris-Edwards}) using a finite difference method~\cite{marenduzzo2007steady}. This method involves approximating the solutions to differential equations by replacing derivatives with algebraic equations. Specifically, our fluid solver uses a five--point stencil to numerically approximate spatial derivatives (up to second order) and verifies accuracy using a predictor--corrector method~\cite{Bhattacharyya2024ActiveNematicsMechanobiology, Li2017NumericalSolutionPDE, Dubois2017LBMFiniteDifference}. The velocity gradients generated by the lattice-Boltzmann solver provide the co--rotational and flow--alignment terms in the Beris-Edwards equation, while the resulting active and elastic stresses are fed back into the Navier-Stokes equations solving for $\mathbf{u}$.

\subsection{Geometry, boundary conditions, and parameters}
\label{sec:geometry_BC_parameters}

While the hybrid lattice-Boltzmann simulation supports three dimensions, the domain used in this paper is a very wide two-dimensional channel of size $L_x = 256$ by $L_y = 256$ with periodic boundaries along $x$ and solid walls at $y=0$ and $y=L_y$. At the walls, no--slip boundary conditions are imposed on the velocity using the bounce-back method, where particles are reflected along their original direction to ensure that the walls are impenetrable and that there is no relative transverse motion between the fluid and boundary~\cite{Krueger2016LatticeBoltzmann}. For the orientational field, we apply free boundary conditions, so there is no imposed preferred alignment at the walls. Any stable near wall orientation therefore comes entirely from the coupled active stress, near-wall shear, and the boundary-layer gradients in $Q$ that develop in confinement.

\begin{wraptable}[10]{r}{0.55\textwidth}
\vspace{-1cm}
\centering
\small
\setlength{\tabcolsep}{4pt}
\renewcommand{\arraystretch}{1.1}
\caption{Simulation parameters used}
\label{table:parameters}
\begin{tabular}{@{}p{0.32\textwidth} l l@{}}
\hline
Parameter & Symbol & Value \\
\hline
Activity & $\zeta$ & $0.0001$--$0.1$ \\
Nematic elasticity & $\kappa$ & $0.0001$--$0.1$ \\
Flow aligning parameter & $\xi$ & $0$--$1.9$ \\
Viscosity & $\eta$ & $0.833$ \\
Bulk free energy scale & $A_0$ & $0.0843$ \\
Reduced temperature & $\gamma$ & $3.0857$ \\
Molecular potential relaxation timescale & $\Gamma$ & {$0.3$} \\
\hline
\end{tabular}
\end{wraptable}

In this work, we are mainly interested in how three main parameters affect active anchoring at solid channel walls: the activity coefficient $\zeta$, the flow–alignment parameter (denoted $\lambda$ in the theory and $\xi$ in the simulation code), and the elastic constant $\kappa$. The rest of the simulation variables (like viscosity, density, temperature, and the equilibrium order parameter) are kept constant across all runs.

While in Sections 2.2 and 2.3 above we use the director-level Leslie--Ericksen parameter $\lambda$, Ludwig uses the Beris--Edwards flow-alignment parameter $\xi$. For a uniaxial nematic with scalar order parameter $S$, the mapping is
\begin{equation}
\lambda = \frac{3S+4}{9S}\,\xi.
\label{eq:lambda_xi_map}
\end{equation}
In our simulations, $S \approx 0.351$ (approximately uniform in the steady state), so $\xi \approx 0.63\,\lambda.$

\pagebreak
The activity $\zeta$ sets the magnitude and sign of the active stress $\boldsymbol{\Pi}_{\mathrm{act}} = -\zeta \mathbf{Q}$ introduced in Section 2. Positive $\zeta$ corresponds to extensile systems that push fluid outwards along their director and draw it in from the sides, while negative $\zeta$ corresponds to contractile systems. We vary $\zeta$ over a range that spans from weakly active states where the bulk flow is relatively ordered, to strongly active states where the interior of the channel exhibits active turbulence. 

The flow–alignment parameter $\lambda$ (or $\xi$) determines whether a nematic is in the flow–tumbling\footnote{Not to be confused with active turbulence, which arises solely from active stresses} regime ($|\lambda|<1$) or the flow–aligning regime ($|\lambda|>1$), and appears in the Leslie–Ericksen equation (Eq. \eqref{eq:Leslie--ericksen}) for the director. In Ludwig the corresponding input parameter is $\xi$, which controls the relative weight of vorticity–induced rotation and strain–induced alignment in the Beris–Edwards dynamics for $\mathbf{Q}$. For fixed values of the other material coefficients, there is a one–to–one mapping between $\xi$ and the effective $\lambda$ that appears in the director limit. Here we vary $\xi$ around the flow–tumbling / flow–aligning crossover and examine how the near–wall anchoring angle changes as we move from a regime where activity dominates to one where passive flow–alignment takes over.

Finally, the elastic constant $\kappa$ sets the cost of spatial gradients in $\mathbf{Q}$, and enters the Landau--de Gennes free energy term in the Beris-Edwards equation. $\kappa$ controls the nematic correlation length and the thickness of boundary layers over which $S$ and $\mathbf{n}$ change near the walls. 

All simulations begin from a uniform nematic state and are run until the system reaches a statistically steady configuration. Observables are then obtained from time--averaged quantities in this steady regime, as explained in the following section.

\subsection{Measured fields and extracted observables}
\label{sec:measured_fields_and_observables}

The Ludwig simulation outputs the velocity field $\mathbf{u}(x,y,t)$, the nematic tensor field $\mathbf{Q}(x,y,t)$, and scalar nematic order parameter field $S(x,y,t)$ at each timestep. We extract the local director $\mathbf{n}$ from the principal eigenvector of $Q$, and compute $\theta$ from $\mathbf{n}$. The equilibrium scalar order parameter $S$ is nearly uniform in space and time, so the simulated nematic orientations can be fully captured by the director field. 

Because this work focuses on active anchoring, the flow field itself is not important for our analysis. Still, it is useful to check that the nematic is exhibiting expected motion in shear flows at the boundaries and, at higher activities, active turbulence in the bulk. Confirming the appearance of shear flows at boundaries is also important for analyzing flow-aligning regimes with $|\lambda|>1$, as the anchoring angle is relative to the direction of shear flow.

We computed wall--angle statistics in the steady--state by pooling samples from both channel walls. We recorded and averaged simulation snapshots every $10000$ timesteps, and for each snapshot, we extracted the near-wall director angle at both walls and combined the resulting samples into a single wall--angle distribution.\footnote{``Near-wall” refers to the first lattice row adjacent to each wall.} This steady--state wall--angle distribution is our primary simulation result, as it shows how the director tends to align relative to the wall for any given regime.  

To confirm that the chosen averaging window lies in steady state, we also plotted time series of the wall--angle mean and mode computed at the same sampling interval, which can be found in Fig.~\ref{fig:steady_state_timeseries} in Appendix B. In these plots, we can clearly observe the preferred orientation angle at walls converge from an initially random state to a steady--state anchored configuration at the channel walls. 

\section{Results}

It is known that active nematic particles at active-passive interfaces tend to align themselves along the interface depending on their activity $\zeta$, with extensile ($\zeta>0$) particles favoring anchoring along the boundary and contractile particles ($\zeta<0$) favoring anchoring perpendicular to the boundary~\cite{blow2014biphasic}. This section delves into the behavior of active nematic particles when instead met with a solid wall.

\subsection{Behavior in the flow-tumbling regime}
\label{sec:behavior_in_flow_tumbling_regime}

We begin by simulating an active nematic with $\xi=0$, in the flow-tumbling regime $|\lambda|<1$. Ludwig uses $\xi$ as the flow-aligning parameter instead of $\lambda$, so with the relation $\lambda = \frac{3S+4}{9S}\,\xi$, the flow-tumbling regime contains flow-aligning parameter values $|\xi| < 0.63$. 

With the averaging process described in Sec.~\ref{sec:measured_fields_and_observables}, we plot the steady--state orientation distribution at the channel walls for extensile and contractile active nematics with $\xi=0$, shown in Fig.~\ref{fig:active-baseline} below. 

\begin{figure}[H]
  \centering
  \includegraphics[width=0.48\textwidth]{combined_angle_dist_tangent_xi_0.000_zeta_0.0020_kappa_0.001.png}
  \includegraphics[width=0.48\textwidth]{combined_angle_dist_tangent_xi_0.000_zeta_-0.0020_kappa_0.001.png}
  \caption{steady--state distribution of the near-wall director angle $\theta$ for (left) extensile and (right) contractile nematic systems ($\zeta=\pm 0.002,\, \kappa=0.001,\, \xi=0$).}
  \label{fig:active-baseline}
\end{figure}

For extensile nematics, the orientation distribution is sharply peaked at $\theta=0^\circ$, corresponding to planar alignment, while contractile nematics align perpendicular to the wall, corresponding to homeotropic alignment. This behavior at solid walls exactly matches that of active anchoring at active-passive boundaries~\cite{blow2014biphasic}.

\subsection{Active anchoring at solid walls}
\label{sec:active_anchoring_at_solid_walls}

To study active anchoring at solid walls, we begin by adapting the derivation for anchoring at active-passive interfaces described in Section 2.3 above. Unlike the active-passive interface case, where the gradient is primarily in $S$, at a solid wall the gradients entering $f^{\rm act}=-\zeta\nabla\cdot Q$ arise from the wall--induced boundary layer in $Q$ produced by confinement and no-slip shear.

We start with the Stokes momentum balance $$\nabla\!\cdot\!\big(2\eta\,\mathbf{E}-\zeta\,\mathbf{Q}\big)=0,
\qquad
\mathbf{E}=\tfrac12\big(\nabla\mathbf{u}+(\nabla\mathbf{u})^{T}\big) = \frac{1}{2}\begin{pmatrix}
    0 & \partial_y u_x \\
    \partial_y u_x & 0
\end{pmatrix}.$$

\pagebreak
We are working in two dimensions, with the $x$-direction along the wall and the $y$-direction being the wall normal. At boundaries, active nematics undergo shear flow $\mathbf{u}(y,t)=(u_x(y,t),0)$ with the no-penetration condition giving $u_y = 0$ at the wall. We assume that fields vary only in the wall normal direction, so $\partial_x(\cdot)=0$. Then, taking the $x$ component of $\nabla\!\cdot\!\boldsymbol{\Pi}=0$ gives
\begin{equation}
    \partial_y\big(2\eta\,E_{xy}-\zeta\,Q_{xy}\big)=0
\end{equation}
This means that the quantity inside the parentheses is independent of $y$, so $$2\eta\,E_{xy}-\zeta\,Q_{xy}=\tau,$$
where $\tau$ is the constant shear stress transmitted across the boundary layer. In an undriven, symmetric channel with no externally applied shear, the time--averaged shear stress vanishes ($\tau=0$), so throughout the boundary layer we have that 
\begin{equation}
    2\eta\,E_{xy}-\zeta\,Q_{xy}=0.
\end{equation}

Finally, using $E_{xy}=\tfrac12\,\partial_y u_x$, we end up with 
\begin{equation}
\eta\,\frac{\partial u_x}{\partial y}
=\zeta\,Q_{xy}.
\end{equation}
In two dimensions, with
$$\mathbf{Q}=2S\big(\mathbf{n}\mathbf{n}-\tfrac12\mathbf{I}\big),
\qquad \mathbf{n}=(\cos\theta,\sin\theta),$$
we have that $Q_{xy}=2S n_x n_y = S\sin 2\theta$, so
\begin{equation}
\eta\,\frac{\partial u_x}{\partial y} = \zeta\,S\,\sin 2\theta.
\label{eq:activity_vel}
\end{equation}

To describe how the director reorients under this shear, we use the Leslie–Ericksen equation~\cite{leslie1968some} for the director field $\mathbf{n}$,
\begin{equation}
    (\partial_t + \mathbf{u}\!\cdot\!\nabla)\mathbf{n}
    - \boldsymbol{\Omega}\!\cdot\!\mathbf{n}
    = \lambda\,\mathbf{E}\!\cdot\!\mathbf{n}
    + \Gamma\,\mathbf{h}_\perp.
    \label{eq:Leslie--ericksen}
\end{equation}

Here $\boldsymbol{\Omega}=\tfrac12\big(\nabla\mathbf{u}-(\nabla\mathbf{u})^T\big)$ is the vorticity tensor, $\lambda$ is the flow–alignment parameter that measures how strongly the director tends to align relative to the flow, $\Gamma$ is a rotational mobility setting the rate at which elastic torques relax distortions in the director field, and $\mathbf{h}_\perp$ is the component of the molecular field perpendicular to $\mathbf{n}$ that drives the director back toward a uniform, low–distortion configuration. The molecular field $\mathbf{h}$ is obtained from the variation of the Landau--de Gennes free energy mentioned above with respect to $\mathbf{n}$ (or $\mathbf{Q}$), and it penalizes splay, bend, and twist of the director~\cite{degennes1993physics,beris1994thermodynamics}.

In this paper's boundary analysis, we can ignore the advective term $(\mathbf{u}\!\cdot\!\nabla)\mathbf{n}$ which represents changes in $\mathbf{n}$ due to transport along the wall. Even though the fluid experiences shear flow close to the wall, we assume that the director is uniform within the boundary layer, so its gradient is zero. Thus, we have that $(\mathbf{u}\!\cdot\!\nabla)\mathbf{n}=0$, so shear affects the director only through the rotation and stretching terms in the Leslie–Ericksen equation, not through advection. For the same reason, we can ignore the molecular field relaxation term at boundaries, as it also depends on spatial variations in $\mathbf{n}$. With these simplifications the director equation reduces to
\begin{equation}
\partial_t \mathbf{n} - \boldsymbol{\Omega}\!\cdot\!\mathbf{n}
= \lambda\,\mathbf{E}\!\cdot\!\mathbf{n}.
\end{equation}
For the shear flow $\mathbf{u}=(u_x(y,t),0)$ introduced above, the nonzero components of $\mathbf{E}$ and $\boldsymbol{\Omega}$ are
$$
\mathbf{E}
=\frac{1}{2}
\begin{pmatrix}
0 & \partial_y u_x\\[4pt]
\partial_y u_x & 0
\end{pmatrix},
\qquad
\boldsymbol{\Omega}
=\frac{1}{2}
\begin{pmatrix}
0 & \partial_y u_x\\[4pt]
-\partial_y u_x & 0
\end{pmatrix}.$$
We can again write the director in terms of the angle $\theta$ measured from the wall,
$$\mathbf{n} =
\begin{pmatrix}
\cos\theta\\[2pt]
\sin\theta
\end{pmatrix},
\qquad
\mathbf{n}_\perp =
\begin{pmatrix}
-\sin\theta\\[2pt]
\cos\theta
\end{pmatrix},$$
where $\mathbf{n}_\perp$ is a unit vector perpendicular to $\mathbf{n}$. The time derivative of $\mathbf{n}$ is
$$\partial_t \mathbf{n}
= \dot\theta
\begin{pmatrix}
-\sin\theta\\[2pt]
\cos\theta
\end{pmatrix}
= \dot\theta\,\mathbf{n}_\perp.$$

We now compute each term in the reduced Leslie–Ericksen equation and then project onto $\mathbf{n}_\perp$ to solve for $\theta(t)$. First, the components of the strain rate tensor $\mathbf{E}$ and the vorticity tensor $\mathbf{\Omega}$ along $\mathbf{n}$ are:

$$\boldsymbol{\Omega}\!\cdot\!\mathbf{n} = \frac{1}{2} \begin{pmatrix} 0 & \partial_y u_x\\[4pt] -\partial_y u_x & 0 \end{pmatrix} \begin{pmatrix} \cos\theta\\[2pt] \sin\theta \end{pmatrix} = \frac{1}{2} \begin{pmatrix} \partial_y u_x \,\sin\theta\\[4pt] -\partial_y u_x \,\cos\theta \end{pmatrix},$$
$$\mathbf{E}\!\cdot\!\mathbf{n} = \frac{1}{2} \begin{pmatrix} 0 & \partial_y u_x\\[4pt] \partial_y u_x & 0 \end{pmatrix} \begin{pmatrix} \cos\theta\\[2pt] \sin\theta \end{pmatrix} = \frac{1}{2} \begin{pmatrix} \partial_y u_x \,\sin\theta\\[4pt] \partial_y u_x \,\cos\theta \end{pmatrix}.$$
We can now project each vector onto $\mathbf{n}_\perp$, to get:
$$\mathbf{n}_\perp\!\cdot\!(\boldsymbol{\Omega}\!\cdot\!\mathbf{n}) =\begin{pmatrix} -\sin\theta & \cos\theta \end{pmatrix} \cdot \frac{1}{2} \begin{pmatrix} \partial_y u_x \,\sin\theta\\[4pt] -\partial_y u_x \,\cos\theta \end{pmatrix} = \frac{1}{2}\,\partial_y u_x \big(-\sin^2\theta - \cos^2\theta\big) = -\frac{1}{2}\,\partial_y u_x,$$
$$\mathbf{n}_\perp\!\cdot\!(\mathbf{E}\!\cdot\!\mathbf{n}) =\begin{pmatrix} -\sin\theta & \cos\theta \end{pmatrix} \cdot \frac{1}{2} \begin{pmatrix} \partial_y u_x \,\sin\theta\\[4pt] \partial_y u_x \,\cos\theta \end{pmatrix} = \frac{1}{2}\,\partial_y u_x\,(\cos^2\theta - \sin^2\theta) = \frac{1}{2}\,\partial_y u_x\,\cos 2\theta.$$
$$\mathbf{n}_\perp\!\cdot\!\partial_t\mathbf{n} = \dot\theta$$
Projecting the full vector equation onto $\mathbf{n_{\perp}}$, we have
$$\mathbf{n}_\perp\!\cdot\!\big(\partial_t\mathbf{n} - \boldsymbol{\Omega}\!\cdot\!\mathbf{n}\big) = \mathbf{n}_\perp\!\cdot\!\big(\lambda\,\mathbf{E}\!\cdot\!\mathbf{n}\big)$$
$$\mathbf{n}_\perp\!\cdot\!\big(\partial_t\mathbf{n} - \boldsymbol{\Omega}\!\cdot\!\mathbf{n}\big)
= \dot\theta - \Big(-\frac{1}{2}\,\partial_y u_x\Big)
= \dot\theta + \frac{1}{2}\,\partial_y u_x$$
$$\mathbf{n}_\perp\!\cdot\!\big(\lambda\,\mathbf{E}\!\cdot\!\mathbf{n}\big) = \lambda\,\mathbf{n}_\perp\!\cdot\!(\mathbf{E}\!\cdot\!\mathbf{n}) = \lambda\,\frac{1}{2}\,\partial_y u_x\,\cos 2\theta.$$
Equating both sides gives
$$\dot\theta + \frac{1}{2}\,\frac{\partial u_x}{\partial y}
= \lambda\,\frac{1}{2}\,\frac{\partial u_x}{\partial y}\,\cos 2\theta,$$
\begin{equation}
\dot\theta
= \frac{1}{2}\,\frac{\partial u_x}{\partial y}\,\big(\lambda\cos 2\theta - 1\big).
\label{eq:theta-shear}
\end{equation}
We can now substitute the active shear relation from Eq. \eqref{eq:activity_vel} above,
$$\frac{\partial u_x}{\partial y} = \frac{\zeta S}{\eta}\,\sin 2\theta,$$
to get a closed evolution equation for the anchoring angle:
\begin{equation}
\dot\theta
= \frac{\zeta S}{2\eta}\,\sin 2\theta\,\big(\lambda\cos 2\theta - 1\big).
\label{eq:theta-active}
\end{equation}
This equation tells us how $\theta$ changes in time due to the combined effect of active stresses (through $\zeta$), the local order $S$, the viscosity $\eta$, and the flow–alignment parameter $\lambda$.\footnote{It also hints at a third steady solution for $\theta$ being $\theta_L = \frac{1}{2}\arccos\!\left(\frac{1}{\lambda}\right)$, introduced in the flow-aligning vs flow-tumbling section below.}

To analyze stability, we look at small perturbations around steady orientations where $\dot\theta=0$. Same as for active anchoring at an active-passive interface, there are two steady states being
$$\sin 2\theta_0=0
\;\;\Rightarrow\;\;
\theta_0=0^\circ \text{(planar anchoring)} \text{ and }
\theta_0=90^\circ \text{(homeotropic anchoring)}.$$

Let
$$\theta(t) = \theta_0 + \delta\theta(t),
\qquad |\delta\theta|\ll1.$$
Expanding to first order in $\delta\theta$, we get that
$$\sin 2\theta = \sin\big(2\theta_0 + 2\delta\theta\big)
\approx \sin 2\theta_0 + 2\cos 2\theta_0\,\delta\theta,$$
$$\cos 2\theta = \cos\big(2\theta_0 + 2\delta\theta\big) \approx \cos 2\theta_0 - 2\sin 2\theta_0\,\delta\theta.$$
Then, as $\sin 2\theta_0=0$ for both planar and homeotropic anchoring, we have
$$\sin 2\theta \approx 2\cos 2\theta_0\,\delta\theta, \quad
\cos 2\theta \approx \cos 2\theta_0.$$
Substituting these into \eqref{eq:theta-active} gives, to first order in $\delta\theta$,
\begin{equation}
    \delta\dot{\theta} = \frac{\zeta S}{4\eta}\,\big(2\cos 2\theta_0\,\delta\theta\big)\,\big(\lambda\cos 2\theta_0 - 1\big) = \frac{\zeta S}{2\eta}\,\cos 2\theta_0\,\big(\lambda\cos 2\theta_0 - 1\big)\,\delta\theta.
\end{equation}
We can write this as $\delta\dot{\theta} = s(\theta_0)\,\delta\theta,$ where
\begin{equation}
    s(\theta_0) = \frac{\zeta S}{2\eta}\,\cos 2\theta_0\,\big(\lambda\cos 2\theta_0 - 1\big).
\end{equation}
is the linear growth rate of small perturbations. 

$\lambda$ is the flow–alignment parameter, which determines whether a nematic aligns with or tumbles in a shear flow. For now, we will focus on the flow-tumbling regime ($|\lambda|<1$), where the only steady state solutions for $\theta$ are $0^\circ$ and $90^\circ$. We will cover the flow-aligning regime ($|\lambda|>1$) where flow-driven orientations emerge in Section~\ref{sec:flow-tumbling-vs-flow-aligning} below.

In the flow-tumbling regime, the factor $(\lambda\cos 2\theta_0 - 1)$ is negative for both $\theta_0=0^\circ$ and $\theta_0=90^\circ$, so the sign of $s(\theta_0)$ is controlled entirely by the factor $\zeta\,\cos 2\theta_0$. From here, we can find the stable and unstable configurations for extensile and contractile systems. 

\pagebreak
For extensile systems $(\zeta>0)$:
$$s(0^\circ)\;\propto\; -\,\zeta\cos 0^\circ = -\,\zeta < 0
\quad \text{(stable)}$$
$$s(90^\circ)\;\propto\; -\,\zeta\cos 180^\circ = \,\zeta > 0
\quad \text{(unstable)}.$$

For contractile systems $(\zeta<0)$:
$$s(0^\circ)\;\propto\; -\,\zeta\cos 0^\circ = -\,\zeta > 0
\quad \text{(unstable)}$$
$$s(90^\circ)\;\propto\; -\,\zeta\cos 180^\circ = \,\zeta < 0 \quad \text{(stable)}.$$

Thus, for extensile systems, planar anchoring is stable while homeotropic anchoring is unstable, while the opposite is true for contractile systems. This result suggests that the stable configuration for active anchoring at solid walls is the same as for nematic particles at active-passive interfaces, confirming our observations in Sec.~\ref{sec:behavior_in_flow_tumbling_regime} above.

\subsection{Flow-tumbling vs flow-aligning regime}
\label{sec:flow-tumbling-vs-flow-aligning}

We have now determined that active nematics at solid walls, in the same way as at active-passive interfaces, undergo active anchoring in flow-tumbling regimes with their stable anchoring configuration dependent on their activity (extensile vs contractile). In this section, we briefly place this result in the broader context of the classical distinction between flow–tumbling and flow–aligning nematics, and then comment on how activity comes in.

To isolate the effect of the flow–alignment parameter $\lambda$, we can consider a simple ``passive" nematic system with no activity ($\zeta = 0$), in an imposed steady shear flow
\begin{equation}
\mathbf{u}(y) = \big(u_x(y), 0\big), \qquad \frac{\partial u_x}{\partial y} = \dot\gamma = \text{const}.
\end{equation}
Our Leslie--Ericksen result (Eq. \eqref{eq:theta-shear}) for the director then reduces to 
\begin{equation}
\dot\theta
= \frac{1}{2}\,\dot\gamma\,\big(\lambda\cos 2\theta - 1\big).
\label{eq:theta-passive}
\end{equation}
Orientations are steady when $\dot\theta = 0$, so from \eqref{eq:theta-passive}, we have
$$
\lambda\cos 2\theta_0 - 1 = 0
\quad\Rightarrow\quad
\cos 2\theta_0 = \frac{1}{\lambda}.
$$
Real solutions exist only if $|1/\lambda|\le 1$, i.e.\ if $|\lambda|\ge 1$. In this regime, we get the classical Leslie angle~\cite{beris1994thermodynamics}
\begin{equation}
\theta_L = \frac{1}{2}\arccos\!\left(\frac{1}{\lambda}\right),
\qquad |\lambda|\ge 1,
\end{equation}
at which the vorticity–induced rotation is exactly balanced by the strain–induced alignment. Thus, for a passive nematic system with $|\lambda|>1$, the director relaxes to a steady angle $\theta_L$ relative to the flow and then remains fixed. For $|\lambda|<1$, the equation $\cos 2\theta_0 = 1/\lambda$ has no real solution, so the system is in the flow-tumbling regime. In this case, the director never settles to a fixed angle, instead continuously rotating or ``tumbling" from the vorticity in the shear flow. 

Returning to active nematics, we have seen that at solid walls, nematic particle activity selects a stable anchoring angle determined by the sign of the activity $\zeta$. In the flow–aligning regime $|\lambda|>1$, however, the shape of nematic particles also drives their director field to align at the Leslie angle $\theta_L$ with respect to the shear. 

In Sec.~\ref{sec:active_anchoring_at_solid_walls} above, in addition to the fixed points at $\theta=0^\circ$ and $\theta=90^\circ$ (from $\sin 2\theta=0$), Eq.~\eqref{eq:theta-active} $$\dot{\theta} \;=\; \frac{\zeta S}{2\eta}\,\sin 2\theta\,(\lambda\cos 2\theta-1)$$ also gives a steady solution at $$\theta_L = \frac12 \arccos\!\left(\frac{1}{\lambda}\right)$$ for $|\lambda|>1$. To determine the stability of $\theta_L$, let $\theta(t)=\theta_L+\delta\theta(t)$ with
$|\delta\theta|\ll 1$. Expanding Eq.~\eqref{eq:theta-active} to first order,
$$\sin 2\theta = \sin\big(2\theta_L + 2\delta\theta\big)
\approx \sin 2\theta_L + 2\cos 2\theta_L\,\delta\theta,$$
$$\cos 2\theta = \cos\big(2\theta_L + 2\delta\theta\big) \approx \cos 2\theta_L - 2\sin 2\theta_L\,\delta\theta.$$
We then have that $$\lambda\cos 2\theta - 1 \approx \lambda(\cos 2\theta_L - 2\sin 2\theta_L\,\delta\theta) - 1
= -2\lambda\sin 2\theta_L\,\delta\theta,$$ with $\lambda\cos 2\theta_L-1=0$. Plugging back into Eq.~\eqref{eq:theta-active}, we have that 
\begin{align}
\dot{\theta} \;&=\; \frac{\zeta S}{2\eta}\,(\sin 2\theta_L + 2\cos 2\theta_L\,\delta\theta)\,(-2\lambda\sin 2\theta_L\,\delta\theta)\notag\\
\;&\approx\; \frac{\zeta S}{2\eta}\,(\sin 2\theta_L)\,(-2\lambda\sin 2\theta_L\,\delta\theta)\notag\\
\;&=\; -\frac{\zeta S}{\eta}\,\lambda\,\sin^2(2\theta_L)\,\delta\theta.
\label{eq:leslie_linearized}
\end{align} to first order in $\delta\theta$. Thus, writing $\delta\dot{\theta} = s_L\,\delta\theta$ gives
\begin{equation}
s_L = -\frac{\zeta S}{\eta}\,\lambda\,\sin^2(2\theta_L).
\label{eq:leslie_growth_rate}
\end{equation} as the linear growth rate of small perturbations about the Leslie angle. 

Because $S>0$, $\eta>0$, and $\sin^2(2\theta_L)\ge 0$, the sign of $s_L$ is set by $-\zeta\lambda$. Thus, when $\lambda>1$, the Leslie angle $\theta_L$ is only stable for extensile activity and unstable for contractile, while for $\lambda<-1$ the opposite is true.

For confined active nematics in the flow-aligning regime, the active torques that favor planar or homeotropic alignment at boundaries compete with the passive tendency to align at $\theta_L$. We test this competition and resulting behavior through numerical simulation in the sections below.



\subsection{Transition to the flow-aligning regime}
\label{sec:transition_to_flow_aligning}

Now, we vary the simulation flow–alignment parameter $\xi$, while keeping $\zeta$ and $\kappa$ constant. For each value of $\xi$, we compute the near-wall director-angle distribution by averaging across both channel walls over the same steady--state window described in Section 3.3 above. Each wall--angle histogram shown below uses 30 uniform bins over the full director angle range, and the wall--angle mode is defined as the peak of the combined distribution.\footnote{Because a histogram-based mode can be sensitive to bin count and bin origin, we checked that varying the number of bins (e.g. 20–60) produces only small changes in the inferred peak location and does not change the qualitative crossover behavior with $\xi$. This can be seen in Fig.~\ref{fig:multibins} in Appendix B.} 

\vspace{-0.3cm}
\begin{figure}[H]
    \centering
    \includegraphics[width=1.0\textwidth]{xi_extensile_filtered_z=0.002.png}
\end{figure}
\vspace{-1cm}
\begin{figure}[H]
    \includegraphics[width=1.0\textwidth]{xi_contractile_filtered.png}
    \caption{Wall director angle distributions for $\xi = 0.6, 0.9, 1.2$ with $\kappa = 0.001$ and $\zeta = \pm0.002$.}
    \label{fig:xi-scan}
\end{figure}

Up to and slightly after the flow–tumbling threshold $\xi=0.63$, the wall anchoring angle is dominated by the active torque generated by the near–wall shear. As $\xi$ increases past this threshold, passive flow–alignment begins to compete with active anchoring, and the steady wall angle shifts away from the purely active configuration and toward the Leslie angle predicted by the passive Leslie–Ericksen model. For extensile active nematics, the Leslie angle configuration is stable, which is reflected as the nematic orientation at the walls almost fully aligns at the Leslie angle relative to the shear flow for higher $\xi$ values. The transition to the flow-aligning regime can be visualized by plotting the mode orientation angle vs flow-aligning parameter, as shown in Fig.~\ref{fig:baseline_leslie} below. For contractile active nematics, as $\xi$ increases the passive flow-alignment still competes with the active anchoring, but the instability of the Leslie angle solution prevents the wall director from settling at $\theta_L$, instead remaining close to the homeotropic active anchoring configuration.

\begin{figure}[H]
    \centering
    \includegraphics[width=0.65\textwidth]{baseline_mode_vs_leslie_xi_final.png}
    \vspace{-0.5cm}
    \caption{Mode of the steady--state wall director angle as a function of the flow-alignment parameter.}
    \label{fig:baseline_leslie}
\end{figure}

Notably, the steady--state wall orientation does not immediately collapse onto the Leslie prediction. Instead, for intermediate values of $\xi$ (approximately $0.63$ through $1.0$), the measured wall--angle mode gradually shifts toward $\theta_L$. This delayed approach is consistent with a competition between two distinct alignment mechanisms, being the near-wall active anchoring driven by boundary-layer active stress, and flow-alignment driven by bulk shear that favors relaxation toward the Leslie angle. As $\xi$ increases, the flow alignment strengthens relative to the active anchoring contribution, and the selected steady--state orientation correspondingly transitions from the active--anchored value towards the Leslie angle. We observe this same transition through a range of different activities and elasticities below, suggesting that it is an inherent feature of active nematics. 

\subsection{Effect of activity strength}
\label{sec:effect_of_activity_strength}

We now study the effect of varying the activity coefficient $\zeta$. The averaged steady--state director angle distribution at the channel walls is shown in Fig.~\ref{fig:zeta-scan} below.
\vspace{-0.2cm}
\begin{figure}[H]
    \centering
    \includegraphics[width=0.95\textwidth]{zeta_extensile_filtered_final.png}
\end{figure}
\vspace{-0.8cm}
\begin{figure}[H]
\centering
    \includegraphics[width=0.95\textwidth]{zeta_contractile_filtered_final.png}
    \vspace*{-0.4cm}
    \caption{Wall director angle distributions for varying values of $\zeta$ with $\xi = 0$ and $\kappa = 0.001$.}
    \label{fig:zeta-scan}
\end{figure}

For $\zeta > 0$ (extensile), the director angle approaches planar alignment, while for $\zeta < 0$ (contractile), the director approaches homeotropic alignment. These trends agree with the predictions from the continuum active anchoring theory. For $|\zeta|$ on the order of $10^{-4}$, the wall angle does not exhibit a clear, persistent peak and fluctuates broadly, indicating that activity is too weak to produce robust anchoring. As $|\zeta|$ increases, the near–wall shear strengthens, and the measured steady--state director angle goes to its active--anchored configuration even as the bulk becomes strongly time-dependent. Finally, at very large activity ($|\zeta|\sim 0.1$), we observe a clear broadening of the near-wall angle distribution compared to intermediate-$|\zeta|$ runs. This is consistent with the bulk entering a strongly turbulent state in which time-dependent shear and defects intermittently perturb the near-wall director away from the anchored configuration. The active anchoring mechanism still selects the same preferred orientation, but the increased amplitude and intermittency of bulk-driven fluctuations reduce the effective stability of the anchored state at the boundary.

As we vary the flow-aligning parameter $\xi$ across nematics with different activities, we observe the same transition from active anchoring to the Leslie angle, as shown in Fig.~\ref{fig:leslie_angle_multiple_zeta} below.

\begin{figure}[H]
    \centering
    \includegraphics[width=0.8\textwidth]{mode_vs_xi_overlay_zeta_with_leslie.png}
    \caption{steady--state wall director angle mode as a function of the flow-alignment parameter $\xi$ for different activity strengths $\zeta$.}
    \label{fig:leslie_angle_multiple_zeta}
\end{figure}

Across the range of activities that exhibit strong active anchoring in the flow-tumbling regime, the transition from active anchoring to the Leslie angle at walls in the steady state is relatively unchanged. We did not include activities greater than $0.01$ or less than $0.001$ because too little and too much activity both disrupt steady state active anchoring, as can be seen in Fig.~\ref{fig:zeta-scan} above.

Within the resolution of our $\xi$ sweep (steps of 0.1) and fluctuations in the mode estimate for the wall angle distribution's peak, both the location of the crossover region and the subsequent approach towards the Leslie angle are similar across intermediate activity strengths. This result indicates that changing $|\zeta|$ primarily modifies the fluctuation amplitude of the wall--angle distribution, and does not shift the $\xi$ scale at which the Leslie--angle mechanism dominates.

\subsection{Effect of elasticity}
\label{sec:effect_of_elasticity}

Finally, we examine the effect of varying the elastic constant $\kappa$. In the continuum picture, the local wall--angle dynamics in Eq. \eqref{eq:theta-active} depend on $\zeta$, $\eta$, $S$, and the flow-alignment parameter $\lambda$, but do not contain $\kappa$ explicitly. Elasticity instead enters by setting how large we expect well--aligned anchoring domains to be, so changing $\kappa$ is expected to primarily affect the width and fluctuation level of the near-wall angle distribution rather than shifting the preferred angle itself.
\vspace{-0.1cm}
\begin{figure}[H]
    \centering
    \includegraphics[width=1.0\textwidth]{kappa_extensile_filtered_final.png}
    \vspace{-0.4cm}
    \caption{Wall director angle distributions for $\kappa$ at varying orders of magnitude with $\zeta = 0.002$ and $\xi = 0$.}
    \label{fig:kappa-scan}
\end{figure}

In the flow-tumbling regime, we find that at low elasticity (around $10^{-4}$), the wall--angle distribution is noticeably broader than the baseline ($\kappa=0.001$), indicating weaker and more fluctuating anchoring. In contrast, for intermediate to large elasticity ($\kappa=10^{-3}$ through $10^{-1}$), the near-wall orientation becomes sharply peaked and remains stable at the expected active anchoring configuration. We can interpret this result as $\kappa$ mainly controlling the strength of the elastic restoring torque and the suppression of near-wall distortions, so at lower $\kappa$ values the wall layer is more easily perturbed by bulk shear and defect activity. At larger $\kappa$ values, in the presence of distortions the elasticity rapidly restores the active anchoring configuration at walls, yielding the expected narrow and stable peak. 

Plotting the mode orientation angle vs flow-aligning parameter $\xi$ over runs with intermediate elasticities ($\kappa=0.001$ to $0.03$), we arrive again at the expected transition of the steady--state wall angle from active anchoring to flow-aligning. 

\begin{figure}[H]
    \centering
    \includegraphics[width=0.65\textwidth]{mode_vs_xi_overlay_kappa_with_leslie.png}
    \caption{steady--state wall director angle mode as a function of the flow-alignment parameter $\xi$ for different elasticities $\kappa$.}
    \label{fig:leslie_angle_multiple_kappa}
\end{figure}

\subsection{Active anchoring penetration depth}
\label{sec:penetration_depth}

Throughout this report, we dropped the elastic contribution to the shear stress at interfaces and walls so that the active and viscous terms give the local relation $\eta\,\partial_y u_x = \zeta\,Q_{xy}$. While this simplification is appropriate for obtaining the shear-driven angle in a thin boundary layer, the penetration depth over which the wall--induced distortion in $\mathbf{Q}$ relaxes into the bulk is expected to be controlled by elasticity, as spatial gradients of $\mathbf{Q}$ cost elastic free energy and generate restoring stresses/torques~\cite{beris1994thermodynamics}.

A standard estimate for the active length scale over which active and elastic stresses balance in active nematics comes from balancing the active stress scale $|\Pi^{\rm act}|\sim|\zeta|\,|Q|$ against the elastic stress scale associated with distortions over a length $\ell$, $|\Pi^{\rm el}|\sim \kappa\,|Q|/\ell^{2}$~\cite{doostmohammadi2018active,Hemingway2016CorrelationLengths}. 
\begin{equation}
|\zeta|\,|Q|\;\sim\;\frac{\kappa\,|Q|}{\ell^{2}}
\qquad\Rightarrow\qquad
\ell_{\rm act}\;\sim\;\sqrt{\frac{\kappa}{|\zeta|}}\,.
\end{equation}

We defined the penetration depth from our simulations as the first lattice row away from the wall where the time and position-averaged director angle approaches the bulk value ($45^\circ$), within a tolerance of $\pm 10^\circ$. Varying the tolerance within a range of ($3^\circ$-$20^\circ$) shifts the extracted threshold rows by only a few lattice spacings and maintains the expected linear scaling of the anchoring length with $\sqrt{\kappa/|\zeta|}$.

\begin{figure}[H]
    \centering
    \includegraphics[width=0.95 \textwidth]{threshold_row_all_datasets.png}
    \caption{Threshold row for wall alignment plotted against $\sqrt{\kappa/|\zeta|}$ for variations in $\zeta$ and $\kappa$.}
    \label{fig:threshold_length}
\end{figure}

We can see a clear correlation between the characteristic penetration depth over which activity imposes its preferred orientation at solid walls and the active anchoring length, $\ell_{\mathrm{act}} \sim \sqrt{\frac{\kappa}{|\zeta|}}.$ 

\section{Conclusion}

In this work, we studied how confined active nematics select preferred director orientations at solid, no-slip walls in the absence of any imposed surface anchoring. Using a continuum description based on Stokes flow with an active stress $\Pi^{\rm act}=-\zeta Q$ and Leslie--Ericksen director dynamics, we derived a closed near-wall angle evolution equation showing that, in the flow-tumbling regime ($|\lambda|<1$), activity alone selects planar anchoring for extensile systems ($\zeta>0$) and homeotropic anchoring for contractile systems ($\zeta<0$). Hybrid lattice--Boltzmann simulations in Ludwig reproduced these predicted stable wall states, as well as showed how increasing the flow-alignment parameter $\xi$ causes a transition from activity-selected anchoring toward a well-defined Leslie angle relative to the wall shear in the flow-aligning regime ($|\lambda|>1$). By varying activity and elasticity, we found that changing $|\zeta|$ primarily affects the sharpness of the wall--angle distribution, while $\kappa$ controls the thickness of the near-wall boundary layer without shifting the preferred angle in the active--anchoring regime. Finally, we confirmed that the active anchoring penetration depth is proportional to the active anchoring length scale $\ell_{\rm act}\sim\sqrt{\kappa/|\zeta|}$.

Going forward, natural next steps would be to extend the solid-wall active anchoring analysis to additional boundary conditions, such as partial/free slip or weak imposed anchoring, and to additional geometries like circular confinement. These results would give further insight into how active nematics behave in the real world, as well as help test the validity of the findings in this report.

\section*{Acknowledgements}

I would like to thank Professor Bialek and Professor Košmrlj for agreeing to be my Princeton advisors and for their guidance and support throughout this project. This experience has been the perfect opportunity to build on the foundational work from my summer internship (much of which consisted of re-deriving theoretical results like Section 2 and Appendix A), and extend it to reach a substantial physical result. I would also like to thank Professor Yeomans for giving me the opportunity to study active nematics with her group, and for her guidance throughout my introduction to the field. Finally, I would like to thank Dr. Jan Rozman, Dr. Ioannis Hadjifrangiskou, and Professor Thampi for their guidance throughout my summer internship, as well as their continued support in my research on active anchoring at solid walls; without their help, I could never have gotten close to where I am now.

\pagebreak
\section*{Appendices}

\subsection*{Appendix A: active stress}

As mentioned in Section 2.2 above, the active stress contribution to the total stress tensor is obtained through the coarse-grained effect of many microscopic force dipoles~\cite{Bhattacharyya2024ActiveNematicsMechanobiology, simha2002hydrodynamic}. Below is a rederivation of this result.

To start, we can consider a dense collection of active nematic particles of length $2\ell$. From here, we can choose any single swimmer $i$ centered at $\mathbf{x}_i$ with unit orientation ${\mathbf{n}}_i$. This swimmer exerts equal and opposite forces $\pm f{\mathbf{n}}_i$ on the fluid at its two ends, which we idealize as the points $\mathbf{x}_i\pm \ell{\mathbf{n}}_i$. The resulting force density produced by particle $i$ is
\begin{equation}
f^{(i)}_\alpha(\mathbf{x})=f n_{i,\alpha}\delta\!\big(\mathbf{x}-\mathbf{x}_i - \ell \mathbf{n}_i\big)-f n_{i,\alpha}\delta\!\big(\mathbf{x}-\mathbf{x}_i+\ell{\mathbf{n}}_i\big),
\label{eq:active_force}
\end{equation}

where the deltas simply mark the two ends of the rod where the forces act. Because $\ell$ is microscopic compared with the length scales on which we are coarse-graining, we can Taylor expand each delta about $\mathbf{x}_i$ to first order to get
$$\delta\!\big(\mathbf{x}-\mathbf{x}_i\mp\ell{\mathbf{n}}_i\big)\approx\delta(\mathbf{x}-\mathbf{x}_i)\mp \ell n_{i,\beta}\partial_\beta \delta(\mathbf{x}-\mathbf{x}_i).$$

Substituting these approximations back into Eq. \eqref{eq:active_force} yields
\begin{equation}
f^{(i)}_\alpha(\mathbf{x})\approx-2f\ell n_{i,\alpha} n_{i,\beta}\partial_\beta \delta(\mathbf{x}-\mathbf{x}_i)=-\partial_\beta\!\left[\sigma_0 n_{i,\alpha} n_{i,\beta}\delta(\mathbf{x}-\mathbf{x}_i)\right],
\end{equation}

where $\sigma_0\equiv 2f\ell$ is the dipole strength of the active nematic. In our low-Re regime, the force density equals the divergence of the stress, $f_\alpha=\partial_\beta \Pi_{\alpha\beta}$. Thus, we have that the single particle contribution to the active stress is
$$\Pi^{(i)}_{\alpha\beta}(\mathbf{x})=-\sigma_0 n_{i,\alpha} n_{i,\beta}\delta(\mathbf{x}-\mathbf{x}_i).$$

To pass into a continuum description, we can then average the contributions of all particles over a small neighborhood around $\textbf{x}$. If we denote $c(\mathbf{x})$ to be the local number density of the active nematic and $\langle  n_\alpha  n_\beta\rangle_{\mathbf{x}}$ to be the orientation average in the neighborhood, coarse-graining then gives an active stress
\begin{equation}
\Pi^{\mathrm{act}}_{\alpha\beta}(\mathbf{x})\approx-\sigma_0c(\mathbf{x})\Big\langle  n_\alpha  n_\beta \Big\rangle_{\mathbf{x}}.
\end{equation}

The second moment $\langle  n_\alpha  n_\beta\rangle$ can be split into an isotropic part and a traceless part,
$$\Big\langle  n_\alpha  n_\beta \Big\rangle_{\mathbf{x}} =\frac{\delta_{\alpha\beta}}{d} +\Big(\Big\langle  n_\alpha  n_\beta \Big\rangle_{\mathbf{x}}-\frac{\delta_{\alpha\beta}}{d}\Big).$$

The isotropic term contributes an equal normal stress in all directions and can be absorbed into a pressure term, as it does not generate shear or couple to orientation. The traceless term is exactly the nematic order tensor (with the simplest normalization)
$$Q_{\alpha\beta}(\mathbf{x}) \equiv \Big\langle  n_\alpha  n_\beta \Big\rangle_{\mathbf{x}}-\frac{\delta_{\alpha\beta}}{d}.$$

Substituting this definition and dropping the isotropic part finally yields the active stress,
\begin{equation}
\Pi^{\mathrm{act}}_{\alpha\beta}(\mathbf{x}) = -\zeta(\mathbf{x})Q_{\alpha\beta}(\mathbf{x}),
\end{equation}
where the activity coefficient $\zeta(\mathbf{x}) \equiv \sigma_0c(\mathbf{x})$ determines the strength of the activity in a given region. Taking $\zeta$ to be constant with approximately uniform particle strength and density, we arrive at $$\boldsymbol{\Pi}_{\mathrm{act}} = -\zeta \textbf{Q}.$$

\subsection*{Appendix B: extra figures}

\begin{figure}[H]
  \centering
  \includegraphics[width=0.48\textwidth]{baseline_mean_wall_angle_vs_time.png}
  \includegraphics[width=0.48\textwidth]{baseline_mode_wall_angle_vs_time.png}
  \caption{Baseline run ($\zeta = 0.002$, $\xi = 0$, $\kappa = 0.001$) (Left) Time series of the mean near-wall angle. (Right) Time series of the near-wall angle mode.}
  \label{fig:steady_state_timeseries}
\end{figure}

\begin{figure}[H]
    \centering
    \includegraphics[width=0.95\textwidth]{xi_extensile_filtered_20bins.png}
    \includegraphics[width=0.95\textwidth]{xi_extensile_filtered_60bins.png}
    \caption{Wall director angle distributions for $\xi = 0.6, 0.9, 1.2$ with $\kappa=0.001$ and $\zeta=0.002$, with bin counts of 20 (above) and 60 (below).}
    \label{fig:multibins}
\end{figure}

\begin{figure}[H]
    \centering
    \includegraphics[width=0.73\textwidth]{q_v_fields_00020000.png}
\end{figure}
\vspace{-0.35cm}
\begin{figure}[H]
    \centering
    \includegraphics[width=0.73\textwidth]{q_v_fields_00050000.png}
\end{figure}
\vspace{-0.35cm}
\begin{figure}[H]
    \centering
    \includegraphics[width=0.73\textwidth]{q_v_fields_00200000.png}
    \caption{Director and velocity fields at various simulation timesteps in a $16\times16$ channel. Steady state emerges at around t = 20000, and we average from t = 100000 to 300000.}
    \label{fig:baseline-orientations}
\end{figure}

Figure \ref{fig:baseline-orientations} above shows that extensile activity generates shear flow near the walls and turbulence in the bulk, along with a constant well-defined anchoring angle after reaching a steady state. This result is consistent with the continuum prediction that extensile systems tend toward planar alignment at boundaries.

\begin{figure}[t]
  \centering
  \includegraphics[width=0.72\textwidth]{combined_wall_order_parameter_mean_vs_time.png}
  \caption{Mean wall order parameter $S$ vs time for the baseline run ($\zeta=0.002, \kappa=0.001, \xi=0$). The near-wall magnitude of nematic order is approximately constant at $S$ = 0.351 over the analysis window.}
  \label{fig:appendixB_wallS_mean_vs_time}
\end{figure}

% ---------- B.2: Wall-order-parameter histograms (3 panels) ----------
\begin{figure}[H]
  \centering
  \includegraphics[width=0.49\textwidth]{wall_S_hist_t=100000.png}\hfill
  \includegraphics[width=0.49\textwidth]{wall_S_hist_t=200000.png}\hfill
  \includegraphics[width=0.49\textwidth]{wall_S_hist_t=300000.png}
  \caption{Distributions of the wall order parameter $S$ at representative times in the baseline run. These histograms confirm that the wall-order distribution remains narrow and stable in time.}
  \label{fig:appendixB_wallS_hists}
\end{figure}

% ---------- B.3: Leslie-angle comparison plots (2 panels) ----------
\begin{figure}[H]
  \centering
  \includegraphics[width=0.49\textwidth]{baseline_mode_vs_leslie_neg_xi_z=-0.002.png}\hfill
  \includegraphics[width=0.49\textwidth]{baseline_mode_vs_leslie_neg_xi_z=-0.02.png}
  \caption{Mode wall angle vs flow-aligning parameter $\xi$ for contractile activity ($\zeta < 0$). These results show that for contractile nematics, the Leslie angle is a stable solution for negative $\xi$ values.}
  \label{fig:appendixB_leslie_comparison}
\end{figure}

\printbibliography

\end{document}